\pg{346} \emph{acto} na \emph{potencia}. Logo o ente, que é \emph{finito}, consta de acto e de \emph{potencia} \footnote{O \emph{ser}, por si, signilica \emph{acto}, \emph{perfeição}, e por isso exclue do seu conceito a \emph{limitação}, que importa imperfeição. Se, pois, o ser se acha limitado, a unict.razão é porque é recebido, como acto na potencia, numa certa essencia, á qual tem uma enpacidade finita. D' ahi a composição do ente finito.}. .

68. O ente infinito é necessario; o \textbf{ente necessario} é infinito.

a) \emph{O ente infinito é necessario}. --- Ente \emph{infinito} é o que encerra todas as perfeições possiveis, e por isso é o proprio ser subsistente, cuja existencia é uma e a mesma coisa com a essencia. Ora um ente, cuja existencia é uma ea mesma coisa com a essencia, existe pela exigencia da sua essencia, não pode ser privado da existencia, sem ser privado da essencia, e por isso é \emph{necessario}. Logo o ente \emph{infinito é necessario}.

b) O \emph{ente necessario é infinito}. --- Ente \emph{necessario} é O que existe pela exigencia da sua natureza, é o proprio ser subsistente, acto purissimo, sem sombra de potencialidade. Ora um ente que é o proprio ser subsistente, acto purissimo, sem sombra de potencialidade, encerra todas as perfeigoes possiveis, e por isso é \emph{infinito} ; pois; se-não-fosse-infinito;-por lhe faltar alguma per- . feição, seria em poólencia para.ella e deixaria de ser acto purissimo. Logo o ente \emph{necessario} é \emph{infinito} ().\footnote{Tambem a verdade d*estas proposições se conhece pela explicação dos lermos. --- \emph{O} ente \emph{infinito} possue todas as perfeições possiveis e por isso exclue toda e qualquer potencialidade, que é imperfeição. Excluindo toda e qualquer potencialidade, \emph{o} ente \emph{infinito} exclue tambem a potencialidade para a não-existencia, e por isso não pode deixar de existir, é absolutamente \emph{necessario}. --- Como tambem o ente \emph{necessario} é existente por si mesmo, pela exigencia da sua essencia, Sendo existente por si mesmo, o ente »\emph{tecessario} é acto purissimo, e porianlo encerra todas as perleições possiveis, é \emph{infinito}. Com eleito, o ente necessario possue o ser pela exigencia da propria essencia. Ora o que um ser possue pela sua essencia não pode ser diminuto ou deficiente. \emph{porque a} essencia não comporta nem o mais nem o menos, mas é \emph{indivisivel}, e por isso o que se possue \emph{pela} exigencia da essencia, possue-se de um modo completo \emph{e} ilimitado. \textsc{Diz-S}. Thomaz: «Quod alicui convenit ex sua natura, el non ex aliqua causa, minoratum in eo el deficiens esse \emph{non} potest. Si enim naturae aliquid essentiale subtrahitur vel additur. \emph{iam altera natura} erit» (\emph{C}. \emph{Gent}, II, 15): Logo 0 ente \emph{necessario} possue o ser \emph{e} a \emph{perfeição} de um modo illimitado e por isso é a plenitude do ente e da \emph{perfeição}, \emph{Quanta} luz irradia dos principios do Angelico Doutor!}

